% 70-18(70쪽) — 물결처럼 오르내리는 $y=f(x)$ 의 그래프(왼쪽 위에서 내려와 $x$축 아래 극소를 지나 극대·극소를 세 번 거치고 오른쪽 위로) · $x=a$, $x=b$ 의 안내 점선. 스캔 화질이 낮아 TikZ 로 다시 그림.
% 원문에 있는 것만: 곡선 · 안내 점선 둘 · 라벨 a, O, b, x, y, y=f(x)(눈금 수치·함수값·c 의 위치는 원문에 없다).
\begin{tikzpicture}[scale=0.45, line width=0.6pt]
  \draw[->, >=stealth, line width=0.7pt] (-3.00,0) -- (5.00,0) node[below=1pt] {$x$};
  \draw[->, >=stealth, line width=0.7pt] (0,-1.15) -- (0,3.75) node[left=1pt] {$y$};
  \draw[densely dotted, line width=0.4pt] (-2.26,0.54) -- (-2.26,0);
  \draw[densely dotted, line width=0.4pt] (4.00,2.39) -- (4.00,0);
  \draw[smooth] plot coordinates {(-2.70,3.15) (-2.58,2.25) (-2.45,1.40) (-2.32,0.72) (-2.20,0.35) (-2.05,-0.10) (-1.90,-0.40) (-1.70,-0.52) (-1.50,-0.42) (-1.32,-0.18) (-1.15,0.15) (-0.95,0.70) (-0.75,1.15) (-0.55,1.48) (-0.35,1.65) (-0.12,1.60) (0.12,1.38) (0.40,1.08) (0.65,0.88) (0.85,0.80) (1.10,0.92) (1.35,1.25) (1.60,1.70) (1.85,2.08) (2.10,2.25) (2.35,2.20) (2.62,1.98) (2.90,1.72) (3.15,1.55) (3.35,1.50) (3.55,1.58) (3.75,1.82) (3.92,2.15) (4.08,2.62) (4.25,3.20)};
  \node[below=1pt, font=\small] at (-2.26,0) {$a$};
  \node[below left=-2pt and -3pt, font=\small] at (0,0) {$\pt{O}$};
  \node[below=1pt, font=\small] at (4.00,0) {$b$};
  \node[font=\small] at (2.95,3.60) {$y=f(x)$};
\end{tikzpicture}
